% sections/methodology.tex
\chapter{Methodology}
\label{chap:methodology}

Questo capitolo descrive il lavoro svolto nell'ambito di questa tesi: viene illustrato come il framework TREAD proposto
da Pallotta et al.~\cite{pallotta} è stato analizzato, re-implementato e (in alcune componenti) esteso e migliorato, per poi
testarlo su un dataset AIS con dati reali relativi all'arcipelago delle Hawaii.

\section{Introduzione}
\label{sec:meth_intro}
Non so bene cosa scrivere qui

\section{Obiettivi}
\label{sec:meth_objectives}

I quesiti di ricerca su cui è stata sviluppata la metodologia sono i seguenti:
\begin{enumerate}
    \item Il framework TREAD (che è stato sviluppato agli inizi degli anni 2010), è ancora applicabile ad un dataset
          moderno e più voluminoso?
    \item Quali aspetti e caratteristiche dell'algoritmo sono ``riciclabili'' e quali hanno richiesto adattamenti o sostituzioni?
    \item È possibile ottenere risultati qualitativamente migliori?
    \item Quali metriche è opportuno utilizzare per quantificare le prestazioni dell'algoritmo?
\end{enumerate} 

\section{Dataset}
\label{sec:meth_dataset}

I dati utilizzati provengono dal portale pubblico del NOAA\footnote{Il NOAA mette a disposizione due modalità di accesso ai dati: tramite il portale \url{https://marinecadastre.gov/accessais/}, che consente di filtrare per zona geografica e intervallo temporale, oppure tramite \url{https://hub.marinecadastre.gov/pages/vesseltraffic}, 
che offre il download in bulk di tutti i dati AIS registrati suddivisi per anno.} (National Oceanic and Atmospheric Administration), un'agenzia scientifica e
normativa statunitense con diversi compiti tra cui il monitoraggio oceanico~\cite{noaa_wiki}.
I file scaricati dalla piattaforma hanno nome nel formato \texttt{AIS\_YYYY\_MM\_DD} (dove YYYY rappresenta l'anno, MM il mese e DD il giorno
di registrazione dei dati) e contengono tutti i messaggi ricevuti dall'ente in una singola giornata di tutto il territorio nazionale monitorato.
Per questa tesi sono stati utilizzati i dati relativi al primo trimestre 2019 (gennaio-marzo, estremi inclusi), filtrati
sull'area delle Hawaii tramite una bounding box le cui estremità hanno le seguenti coordinate:
\begin{itemize}
    \item Latitudine: da $18.77^\circ$N a $22.63^\circ$N,
    \item Longitudine: da $160.11^\circ$W a $154.39^\circ$W.
\end{itemize}

\subsection{Struttura del dataset}
\label{ssec:meth_dataset_structure}

Ogni record del dataset corrisponde a un singolo messaggio AIS trasmesso da una nave e contiene diversi campi, riportati nella tabella~\ref{tab:dataset_columns}.
I campi si dividono in due categorie: attributi statici, che descrivono le caratteristiche fisse della nave e normalmente non variano tra
un messaggio e l'altro, e attributi dinamici, che riflettono lo stato corrente della nave al momento della trasmissione.

\begin{table}[ht]
\centering
\caption{Struttura del dataset AIS del NOAA.}
\label{tab:dataset_columns}
\begin{tabular}{llp{7cm}}
\hline
\textbf{Colonna} & \textbf{Tipo} & \textbf{Descrizione} \\
\hline
\multicolumn{3}{l}{\textit{Attributi statici}} \\
\texttt{MMSI}          & intero   & Maritime Mobile Service Identity: identificativo univoco della nave \\
\texttt{VesselName}    & stringa  & Nome della nave \\
\texttt{IMO}           & stringa  & Numero IMO (International Maritime Organization) \\
\texttt{CallSign}      & stringa  & Nominativo radio della nave \\
\texttt{VesselType}    & intero   & Codice numerico che indica la tipologia di nave (cargo, tanker, passeggeri, ecc.) \\
\texttt{Length}        & intero   & Lunghezza della nave in metri \\
\texttt{Width}         & intero   & Larghezza della nave in metri \\
\texttt{Draft}         & decimale & Pescaggio della nave in decimi di metro \\
\texttt{Cargo}         & intero   & Codice del tipo di carico trasportato \\
\texttt{TransceiverClass} & booleano & Classe del transponder AIS: \texttt{A}  o \texttt{B} (imbarcazioni minori) \\
\hline
\multicolumn{3}{l}{\textit{Attributi dinamici}} \\
\texttt{BaseDateTime}  & datetime & Timestamp del messaggio (UTC) \\
\texttt{LAT}           & decimale & Latitudine della posizione (gradi decimali) \\
\texttt{LON}           & decimale & Longitudine della posizione (gradi decimali) \\
\texttt{SOG}           & decimale & Speed Over Ground: velocità rispetto al fondale in nodi \\
\texttt{COG}           & decimale & Course Over Ground: direzione di movimento in gradi (0--360°) \\
\texttt{Heading}       & decimale & Prua della nave in gradi (0--360°), distinta dal COG \\
\texttt{Status}        & intero   & Stato di navigazione AIS (es. in navigazione, all'ancora, in manovra) \\
\hline
\end{tabular}
\end{table}

Ogni record del dataset, contenendo una coppia di coordinate geografiche (\texttt{LAT}, \texttt{LON}),
è rappresentabile come un punto su una mappa. A ciascun punto è inoltre associato un vettore di movimento:
il modulo corrisponde alla velocità istantanea della nave (\texttt{SOG}), mentre la direzione è determinata
dall'angolo \texttt{COG}, espresso in gradi rispetto al nord geografico. Questa rappresentazione --- punto
più vettore --- è alla base delle elaborazioni descritte nelle sezioni successive, ed è il motivo per cui
nel resto del capitolo i termini \emph{record} e \emph{punto} vengono usati in modo intercambiabile.


\subsection{Qualità dei dati}
\label{ssec:meth_dataset_quality}

Oltre alla struttura, è opportuno considerare la qualità intrinseca dei dati. Un vantaggio rilevante del dataset NOAA
rispetto ad altre fonti pubbliche (in particolare i dati danesi considerati e poi scartati in fase esplorativa) è che
i dati sono già stati filtrati dagli errori più grossolani, rendendo l'uso di tecniche di pulizia aggiuntive meno
critico. È stata tuttavia mantenuta l'assunzione esplicita che le colonne \texttt{MMSI} e \texttt{BaseDateTime} 
siano affidabili e prive di errori, in quanto queste quattro colonne costituiscono la base di tutte le elaborazioni successive.

% -----------------------------------------------------------------------
\section{Pre-processing dei dati}
\label{sec:meth_preprocessing}

L'intera fase di pre-processing è implementata in Python, utilizzando dei notebook Jupyter.
Le tecnologie utilizzate verranno descritte meglio nella Sezione~\ref{sec:meth_technologies}.

\subsection{Importazione, filtraggio geografico e rimozione dei duplicati}

I file scaricati dalla piattaforma del NOAA sono in formato ``.zip''. Pertanto vengono automaticamente decompressi, importati in memoria
e successivamente eliminati dal disco per liberare spazio.
Durante l'importazione, ogni file csv letto viene salvato in memoria sottoforma di \texttt{DataFrame}. In questa fase, 
vengono immediatamente scartati tutti quei messaggi le cui coordinate di latitudine e longitudine ricadono
al di fuori del perimetro citato nella sezione~\ref{sec:meth_dataset}.
Tutti i \texttt{DataFrame} così creati vengono concatenati in un unico \texttt{DataFrame} finale, che viene esportato in un unico file CSV di backup.
Le operazioni di importazione e filtraggio avvengono grazie alla libreria \textbf{Polars}, che è stata scelta per il fatto che (a differenza
di Pandas) supporta il multithreading nativamente. 

In una fase successiva, i record esattamente duplicati (identici in ogni attributo) vengono rimossi, lasciandone solo uno per ogni gruppo duplicato.

Viene poi rimosso un ulteriore sottoinsieme di record: quelli generati da navi, la cui somma di record presenti nel database è inferiore
al 10\% della media calcolata sul totale delle navi presenti. 
Per identificare l'appartenenza di un record ad una nave o meno, è stato utilizzato il valore di \texttt{MMSI}. Come citato in precedenza alla
sezione~\ref{ssec:meth_dataset_quality}, tale attributo si assume corretto e validato dal NOAA in fase di raccolta dei dati.

La soglia usata per rimuovere le navi meno attive è calcolata dinamicamente (e non come valore fisso) in modo
da rendere il codice indipendente dalle dimensioni del dataset che si decide di utilizzare. L'idea dietro questa decisione è che
le navi al di sotto di questa soglia contribuiscono in maniera trascurabile al traffico dell'area e la loro presenza potrebbe 
distorcere la fase di clustering.
Prima dell'applicazione di questa soglia il dataset conteneva 602 imbarcazioni distinte (5.460.800 record). In seguito il numero di 
navi distinte è passato a 403 (5.415.168 righe), eliminando circa lo 0,84\% dei record totali. 

\subsection{Calcolo della velocità media e classificazione degli stati}
\label{ssec:meth_status}
 
Una volta completata la fase di importazione e pulizia, l'analisi viene condotta su una nave alla volta utilizzando
il valore di \texttt{MMSI} come chiave primaria.
Per ciascuna nave viene costruito un \texttt{DataFrame} separato, ordinato cronologicamente per \texttt{BaseDateTime} e contenente
tutti messaggi che quell'imbarcazione ha generato nel periodo di analisi all'interno dell'area geografica analizzata.
Per ogni coppia di punti consecutivi appartenenti alla stessa nave viene calcolata la velocità media \texttt{Avg\_Speed},
come rapporto tra la distanza percorsa e l'intervallo di tempo trascorso. La distanza viene calcolata sfruttando la formula
della \emph{Great-Circle Distance}, che tiene conto della curvatura terrestre usando come raggio della Terra il raggio equatoriale definitio
dallo standard WGS-84 e pari a 6\,378\,137 m. La velocità risultante viene convertita in nodi nautici.
 
Il calcolo viene gestito in modo vettorizzato: ogni traiettoria viene rappresentata come due sequenze di punti sfalsate
di un passo, ottenendo le coppie (punto precedente, punto corrente). Nei casi in cui l'intervallo di tempo
tra due punti consecutivi supera la soglia \texttt{LOST\_TIME\_THRESHOLD\_SECONDS} (impostata a 620 s e corrispondente
a circa 10 volte il periodo tipico di trasmissione di un dispositivo AIS di Classe A), la velocità media calcolata
viene scartata e sostituita con il valore di \texttt{SOG}, a condizione che quest'ultimo sia plausibile (inferiore
alla soglia \texttt{ABSURD\_SPEED\_THRESHOLD} di 102 nodi). In caso contrario, la velocità viene ricalcolata
ancora come rapporto distanza/tempo. Al primo punto di ogni traiettoria, per il quale non esiste un punto
precedente, viene assegnato direttamente il valore di \texttt{SOG}. I punti con \texttt{Avg\_Speed} superiore
a 102 nodi vengono eliminati, in quanto ritenuti il risultato di errori di posizionamento GPS.
 
A partire da \texttt{Avg\_Speed}, viene costruita la colonna \texttt{EstimatedStatus}, che classifica ogni punto
in una delle sei categorie seguenti. Le prime quattro riguardano la gestione dei gap temporali e vengono
assegnate durante il calcolo della velocità:
 
\begin{itemize}
    \item \textbf{lost}: punto immediatamente precedente un gap temporale superiore a \texttt{LOST\_TIME\_THRESHOLD\_SECONDS}.
          Indica l'ultimo momento in cui la nave era tracciata prima di scomparire temporaneamente dal sistema.
    \item \textbf{return}: primo punto successivo a tale gap. Indica il momento in cui la nave torna visibile
          dopo una breve interruzione. I punti \emph{lost} e \emph{return} non entrano nella fase di clustering,
          ma sono utili per la corretta interpretazione delle traiettorie.
    \item \textbf{exit}: punto immediatamente precedente un gap superiore a \texttt{EXIT\_TIME\_THRESHOLD\_HOURS}
          (12 ore). Indica che la nave ha lasciato l'universo di analisi o ha smesso di trasmettere per un periodo
          prolungato.
    \item \textbf{entry}: primo punto successivo a tale gap. Indica il rientro della nave nell'universo di analisi
          dopo un'assenza significativa.
\end{itemize}
 
Gli status \emph{exit} e \emph{entry} vengono assegnati dopo \emph{lost} e \emph{return}, sovrascrivendoli quando
il gap supera la soglia più lunga.
Le restanti due categorie vengono assegnate in una seconda fase, sui punti che non hanno ancora ricevuto uno status:
 
\begin{itemize}
    \item \textbf{stationary}: punti con \texttt{Avg\_Speed} $\leq$ \texttt{SPEED\_THRESHOLD} (0.5 nodi).
    \item \textbf{sailing}: tutti i punti rimanenti, in navigazione ordinaria.
\end{itemize}
 
Al termine, il primo punto di ciascuna traiettoria viene forzato a \emph{entry} e l'ultimo a \emph{exit},
indipendentemente dallo status calcolato. Quando più punti consecutivi risultano tutti \emph{entry} o tutti
\emph{exit} (situazione possibile in presenza di gap ravvicinati) viene conservato solo l'ultimo punto
\emph{entry} della sequenza e il primo punto \emph{exit}, eliminando i rimanenti.

%-----------------------------------------------------------------------------------------------------------
\section{Implementazione del framework TREAD e differenze rispetto all'originale}
\label{sec:meth_algorithm}

\subsection{Architettura generale}

Il framework TREAD di Pallotta et al.~\cite{pallotta} si articola in due macro fasi: \emph{Knowledge Discovery} e \emph{Knowledge Exploitation}.
Nella prima fase, il flusso di messaggi AIS viene elaborato per costruire un modello statistico del traffico normale, articolato in tre classi
di oggetti: vessel objects ($Vs$), waypoint objects ($WPS$) e route objects ($Rs$).
Nella seconda fase, le rotte apprese vengono utilizzate per classificare nuove traittorie e individuare comportamenti anomali.

La presente implementazione riproduce questa struttura a due fasi. Tuttavia, il processo non è incrementale in senso stretto:
anzichè aggiornare il modello messaggio per messaggio come nell'algoritmo originale, l'intero dataset viene caricato in memoria
e processato in batch. Questa scelta è stata motivata dai vincoli hardware descritti nella sezione~\ref{sec:meth_limits}.


\subsection{Identificazione dei waypoint: DBSCAN e H-DBSCAN}
\label{ssec:meth_clustering}

Il passo centrale della fase di Knowledge Discovery è l'identificazione dei waypoint. In questo
contesto, un waypoint è un'area geografica appartenente all'universo di analisi in cui le navi:
\begin{itemize}
    \item entrano nell'universo di analisi;
    \item escono dall'universo di analisi;
    \item rallentano significativamente, con velocità media inferiore alla soglia
          \texttt{SPEED\_THRESHOLD} (0.5 nodi), al di sotto della quale la nave è ritenuta
          sostanzialmente stazionaria.
\end{itemize}
A differenza di quanto avviene su terra, in mare è molto difficile per una nave
rimanere completamente ferma per via di correnti, vento e onde che spesso causano qualche tipo di movimento. 
Per questo motivo, anziché cercare punti di velocità nulla, si è deciso di introdurre una soglia
(\texttt{SPEED\_THRESHOLD}) sotto la quale un imbarcazione è considerata stazionaria. 
I waypoint non hanno dimensione fissa: ciascuno occupa una porzione dello spazio geografico 
determinata da come si distribuiscono e da quanti sono i punti che lo compongono. Nella pratica, la maggior parte di essi 
tende ad occupare un'area molto ridotta rispetto all'universo di analisi.
Nell'algoritmo originale l'identificazione dei waypoint viene effettuata tramite un approccio basato su DBSCAN.
La presente implementazione applica prima DBSCAN e poi H-DBSCAN come variante
migliorativa (\emph{Hierarchical DBSCAN}).

\subsubsection{Preparazione degli input per il clustering}
 
L'input fornito agli algoritmi di clustering è costituito da tutti i punti classificati come 
\emph{stationary}, \emph{entry} o \emph{exit}, che rappresentano circa il 65,3\% del dataset totale.

Per ridurre la dimensione dell'input e superare le limitazioni di memoria dell'implementazione scikit-learn di DBSCAN, 
il cui costo di memoria è $O(n^2)$, i punti vengono aggregati tramite un'operazione di \emph{group-by} sulle sole coordinate GPS, 
contando quante volte ogni coppia (\texttt{LAT}, \texttt{LON}) compare nel dataset. Questa operazione è il motivo per cui si è deciso
di arrotondare da 5 a 4 cifre decimali i valori di latitudine e longitudine (vedi Tabella~\ref{tab:dtype_optimization}). così
facendo si riduce il numero di coppie distinte in input all'algoritmo di clustering.
Il conteggio risultante viene interpretato come un \emph{peso} di ciascuna coordinata. 
Per attenuare lo sbilanciamento tra coordinate molto visitate e coordinate rare, i pesi vengono normalizzati applicando 
il logaritmo in base 2; le coordinate con peso zero dopo la trasformazione logaritmica vengono riportate artificialmente a 1.


\subsubsection{DBSCAN}
 
L'algoritmo DBSCAN viene eseguito con i parametri $\varepsilon = 0.075$ e $N_{min} = 10$, scelti empiricamente 
in modo da produrre cluster geograficamente coerenti con i porti e i punti di passaggio dell'arcipelago delle Hawaii. 
Questo passo riproduce l'approccio di Pallotta et al.\ ed è incluso come riferimento per il confronto con la variante H-DBSCAN.

\subsubsection{H-DBSCAN}
 
Il motivo principale per introdurre H-DBSCAN è la sua capacità di identificare cluster di densità variabile, cosa che DBSCAN non 
riesce a gestire bene \cite{schubert2017}. 
Il traffico nelle Hawaii presenta una densità fortemente eterogenea: i porti principali come Honolulu hanno concentrazioni 
di punti molto più elevate rispetto a zone di passaggio o porti minori. 
In questo contesto, i parametri fissi di DBSCAN tendono a sovra-segmentare le aree sparse o a fondere artificialmente cluster distinti 
nelle aree dense.
 
Poiché la libreria \texttt{hdbscan}\footnote{Libreria esterna \texttt{hdbscan} (\url{https://github.com/scikit-learn-contrib/hdbscan}), 
distinta dall'implementazione inclusa in scikit-learn.} non consente di associare i pesi ad ogni punto ma li considera tutti equiponderati, 
la tabella aggregata viene espansa: ogni coordinata viene replicata un numero di volte pari al suo peso logaritmico. 
La tabella espansa conta circa 68.000 righe, una dimensione che consente l'esecuzione in poche decine di secondi. 
H-DBSCAN viene eseguito con i parametri $\text{min\_cluster\_size} = 30$, $\text{max\_cluster\_size} = 10\,000$ 
e $\text{cluster\_selection\_epsilon} = 0.01$.

\subsection{Costruzione dei waypoint e identificazione delle rotte}
\label{ssec:meth_routes}
 
Una volta determinati i cluster tramite H-DBSCAN, per ciascun cluster viene calcolato il \emph{convex hull} (inviluppo convesso) 
tramite \texttt{scipy.spatial.ConvexHull}. L'inviluppo convesso di ogni cluster definisce il perimetro del waypoint corrispondente, 
rappresentato come un poligono tramite la libreria Shapely.
 
Per ogni punto del dataset viene quindi verificato se esso ricade all'interno di uno dei poligoni waypoint, 
informazione salvata nella colonna \texttt{IsInWP}. 
Un waypoint è considerato \emph{significativo} per un dato punto se quest'ultimo soddisfa almeno una delle condizioni seguenti: 
\texttt{Avg\_Speed} $\leq$ \texttt{SPEED\_THRESHOLD}, oppure il suo \texttt{EstimatedStatus} è \emph{entry} o \emph{exit}.


L'assegnazione delle rotte avviene una nave alla volta, all'interno di ogni coppia
(\emph{entry}, \emph{exit}) identificata nel dataset. A ogni punto viene assegnata
un'etichetta secondo la seguente logica:
\begin{itemize}
    \item se il punto si trova all'interno di un waypoint significativo, l'etichetta
          corrisponde al nome di quel waypoint;
    \item se il punto si trova in navigazione tra due waypoint, l'etichetta assume il
          formato \texttt{WP\_X -- WP\_Y}, dove \texttt{WP\_X} e \texttt{WP\_Y} sono
          rispettivamente il waypoint di provenienza e quello di destinazione; nei casi
          in cui uno o entrambi non siano identificabili, viene utilizzato il
          segnaposto \texttt{<NA>}.
\end{itemize}

Con queste ultime operazioni, termina la fase di Knowledge Discovery: ogni punto
del dataset ha associata un'etichetta di rotta secondo i criteri appena descritti, e il
modello di traffico appreso è pronto per essere utilizzato nella fase di Knowledge
Exploitation.

% -----------------------------------------------------------------------
\section{Anomaly detection}
\label{sec:meth_anomaly}
 
Prima di descrivere la fase di anomaly detection, è opportuno introdurre due concetti
fondamentali. Una \textbf{rotta} è un'entità astratta che modella il fatto che almeno una volta una 
nave ha viaggiato dal waypoint di partenza a quello di arrivo (e non il viceversa).
Ipotizzando di modellare i waypoint come nodi di un grafo orientato, una rotta sarebbe uno degli archi.
Un \textbf{trip} indica invece il singolo viaggio compiuto da una specifica nave lungo una rotta. 
A ogni rotta possono quindi essere associati più trip, tanti quante sono le navi che l'hanno percorsa nel
periodo analizzato.

La fase di anomaly detection segue le equazioni del framework TREAD originale \cite{pallotta2013}.
 L'analisi si concentra su una singola rotta selezionata (una delle più trafficate tra quelle estratte nella fase precedente) 
 e utilizza i trip associati a tale rotta come unità di analisi
 
\subsection{Decomposizione della velocità e preparazione delle feature}
 
Per ogni punto del dataset vengono calcolate le componenti cartesiane della velocità $V_x$
e $V_y$ a partire da \texttt{SOG} e \texttt{COG}. Il COG, espresso in gradi rispetto al
nord geografico (angolo \emph{bearing}), viene convertito in angolo aritmetico standard
tramite la formula $\theta = (450 - \text{COG}) \mod 360$, e poi in radianti. Le componenti
risultanti sono:
\[
V_x = \text{SOG} \cdot \cos(\theta), \qquad V_y = \text{SOG} \cdot \sin(\theta)
\]
Pallotta et al.\ \cite{pallotta2013} propongono un approccio alternativo per il calcolo
delle stesse componenti, basato sulla funzione arcotangente e sul teorema di Pitagora.
Questo approccio presenta alcune criticità. In primo luogo, la funzione arcotangente è
definita a meno di un'ambiguità di segno: per ricavare il quadrante corretto della
direzione è necessario gestire esplicitamente i casi in cui $V_x$ e $V_y$ assumono segno
negativo, introducendo logica condizionale aggiuntiva. In secondo luogo, l'uso combinato
di arcotangente e radice quadrata è computazionalmente più oneroso rispetto al calcolo
diretto tramite seno e coseno. Infine, l'approccio inverso introduce potenziali errori
numerici nelle situazioni limite, come velocità molto prossime a zero, in cui il rapporto
tra le componenti diventa instabile. Per questi motivi, nella presente implementazione si
è preferito calcolare $V_x$ e $V_y$ direttamente tramite seno e coseno dell'angolo
$\theta$, ottenendo un codice più semplice ed efficiente.



\begin{table}[ht]
\centering
\caption{Ottimizzazione dei tipi di dato: conversioni effettuate e specifiche tecniche indicate dalla normativa.}
\label{tab:dtype_optimization}
\begin{tabular}{lllp{6.5cm}}
\hline
\textbf{Colonna} & \textbf{Formato originale} & \textbf{Nuovo formato} & \textbf{Motivazione} \\
\hline
\texttt{MMSI} & \texttt{int64} & \texttt{int32} & MMSI è codificato su 30 bit, con valori validi fino a 999\,999\,999. Un \texttt{int32} (massimo $\approx 2.1 \times 10^9$) è sufficiente. \\
\texttt{LAT} & \texttt{float64} & \texttt{float32} & La latitudine è codificata su 27 bit in unità di $1/10\,000$ di minuto, corrispondenti a una risoluzione di circa 0.185 m. La precisione di \texttt{float32} ($\approx 7$ cifre decimali) è ampiamente sufficiente. \\
\texttt{LON} & \texttt{float64} & \texttt{float32} & La longitudine è codificata su 28 bit con la stessa risoluzione della latitudine. Stessa giustificazione. \\
\texttt{SOG} & \texttt{float64} & \texttt{float32} & La SOG è codificata su 10 bit in passi di $1/10$ di nodo, con massimo 102.2 nodi. Nessuna necessità di doppia precisione. \\
\texttt{COG} & \texttt{float64} & \texttt{float32} & Il COG è codificato su 12 bit in passi di $1/10$ di grado (0--359.9°). Stessa giustificazione della SOG. \\
\texttt{Heading} & \texttt{float64} & \texttt{float32} & La prua è codificata su 9 bit in gradi interi (0--359). Stessa giustificazione. \\
\texttt{VesselType} & \texttt{float64} & \texttt{int8} & Il codice tipo nave è codificato su 8 bit (valori 0--255, Tab.~50). Un \texttt{int8} non con segno sarebbe sufficiente; si usa \texttt{int8} con segno per compatibilità con Pandas. \\
\texttt{Status} & \texttt{float64} & \texttt{int8} & Lo stato di navigazione è codificato su 4 bit (valori 0--15, Tab.~46). \\
\texttt{Length} & \texttt{float64} & \texttt{int16} & Le dimensioni nave sono codificate complessivamente su 30 bit suddivisi in quattro campi A, B, C, D da massimo 511/511/63/63 m rispettivamente (Tab.~50). Un \texttt{int16} (massimo 32\,767) è sufficiente. \\
\texttt{Width} & \texttt{float64} & \texttt{int16} & Stessa giustificazione di \texttt{Length}. \\
\texttt{Draft} & \texttt{float64} & \texttt{float32} & Il pescaggio massimo è codificato su 8 bit in passi di $1/10$ di metro, con massimo 25.5 m (Tab.~50). Nessuna necessità di doppia precisione. \\
\texttt{Cargo} & \texttt{float64} & \texttt{int8} & Il codice carico condivide la stessa codifica a 8 bit del tipo nave (Tab.~50). \\
\texttt{TransceiverClass} & \texttt{object} & \texttt{bool} & Il campo assume esclusivamente i valori \texttt{A} (dispositivi SOLAS, Classe A) o \texttt{B} (imbarcazioni minori). Rinominato in \texttt{IsClassA} e convertito in booleano. \\
\texttt{BaseDateTime} & \texttt{object} & \texttt{datetime64} & Conversione da stringa a tipo temporale nativo di Pandas per consentire operazioni aritmetiche su timestamp. \\
\hline
\end{tabular}
\end{table}





Per ognuno dei punti seguenti prevedo una sezione dedicata

- Spiego qual'è l'obiettivo (vedere se l'algoritmo di pallotta è ancora applicabile e se si può migliorare)

- Ripercorro passo passo quello che ho fatto e lo spiego, indicando le differenze con l'algoritmo originale

- Indico le tecnologie utilizzate e perchè le ho scelte

- Valutazione delle performance dell'algoritmo (metriche da stabilire)

- Limiti della metodologia (es: dimensione massima dataset, hardware disponibile, problemi a raccogliere i dati, ecc...)

- Breve conclusione con anticipo del capitolo successivo